\section{Model Architecture and Module Boundaries}
\label{sec:architecture}

\subsection{A) Central emotion-classification model}

For our model architecture we chose a modified ResNet18 Model. Instead of the traditional \( 224 \times 224 \) image resolution our model accepts \( 64 \times 64 \) images and to accommodate for that we made several architectural modifications. One of the main modifications was changing the initial convolution layer, as we replaced ResNet18's \( 7 \times 7 \) kernel with a \( 3 \times 3 \) one to ensure that we do not lose critical information given the lower resolution input. The network then consists of \( 4 \) layers with each \( 2 \) blocks implementing residual stages. Each of the blocks use batch normalization followed by SiLU instead of ResNet's classic ReLU, as the Swish activation function offers smoother gradients and slightly better optimization in deeper networks. 

In addition, after the blocks we use a Global Average Pooling (GAP) instead of the traditional Flatten layer to reduce the total parameter count. Then we implemented a \( 6 \)-node fully connected layer with a softmax function to finish our model off.

\begin{figure}[t]
    \centering
    \includegraphics[width=\linewidth]{media/Architecture-Structure.jpeg}
    \caption{Pipeline overview of the classifier architecture. Starting from a \(64\times64\) input image, the model applies an initial \(3\times3\) convolution, four residual stages with stage-specific strides and feature-map dimensions, Global Average Pooling, and a final 6-node fully connected layer.}
    \label{fig:architecture_pipeline}
\end{figure}

\paragraph{Optional internal branches of the central model}

\textbf{Raw landmark fusion.}
Motivation: inject explicit facial geometry with minimal parameter overhead.
Technical description: for a batch size \(B\), landmarks are represented as \(L\in\mathbb{R}^{B\times 2N}\) and concatenated directly with pooled backbone features \(z_{\text{cnn}}\in\mathbb{R}^{B\times C}\), where \(C\) is the width-multiplied final channel dimension.
No intermediate MLP is applied to \(L\), so coordinate precision is preserved and branch latency remains negligible.
Integration: late fusion at head input, \(z=[z_{\text{cnn}};L]\).

\textbf{Multi-view heatmap branch.}
Motivation: represent landmark information at multiple spatial scales.
Technical description: landmarks are reshaped to point sets \((x_i,y_i)\), clamped to the valid image grid, and rasterized into Gaussian maps for \(\sigma\in\{1,2,4\}\):
\[
H_k(x,y)=\sum_{i=1}^{N}\exp\left(-\frac{(x-x_i)^2+(y-y_i)^2}{2\sigma_k^2+\epsilon}\right).
\]
The stacked tensor \(H\in\mathbb{R}^{B\times 3\times 64\times 64}\) is processed by a compact Conv-BN-SiLU encoder and global average pooling to produce a dense branch embedding \(z_{\text{hm}}\in\mathbb{R}^{B\times d_{\text{hm}}}\).
Integration: \(z_{\text{hm}}\) is concatenated into the classifier head vector.

\textbf{Landmark-guided attention.}
Motivation: amplify expression-relevant regions before deep feature extraction.
Technical description: image and heatmaps are channel-wise concatenated, \(F=[I;H]\in\mathbb{R}^{B\times (1+3)\times 64\times 64}\), then passed through a Conv-BN-SiLU-Conv-Sigmoid attention block to estimate \(A\in\mathbb{R}^{B\times 1\times 64\times 64}\).
Input modulation is performed as
\[
\tilde{I}(x,y)=I(x,y)\cdot(1+A(x,y)).
\]
In parallel, an auxiliary feature path applies Conv-BN-SiLU-\(1\times1\) Conv and adaptive pooling on \(F\), yielding \(z_{\text{att}}\in\mathbb{R}^{B\times d_{\text{att}}}\) (\(d_{\text{att}}=32\) in the current setup).
Integration: \(\tilde{I}\) replaces the original image in the backbone stream; \(z_{\text{att}}\) is appended at head fusion.

\textbf{Sobel feature branch.}
Motivation: preserve high-frequency edge structure as a complementary cue.
Technical description: normalized grayscale input \(I\in[-1,1]\) is remapped to \([0,1]\), convolved with fixed Sobel kernels to obtain \(G_x,G_y\), and converted to gradient magnitude
\[
M=\sqrt{G_x^2+G_y^2+\epsilon}.
\]
\(M\) is then normalized per sample by \(\max_{x,y} M(x,y)\) for contrast-invariant scaling and encoded by a dedicated Conv-BN-SiLU branch with global pooling, producing \(z_{\text{sobel}}\in\mathbb{R}^{B\times d_{\text{sobel}}}\).
Integration: \(z_{\text{sobel}}\) is fused by concatenation at the classifier head.

\textbf{Landmark GNN branch.}
Motivation: capture relational geometry between landmarks, not only absolute coordinates.
Technical description: landmarks are interpreted as \(N\) nodes with 2D features, \(V\in\mathbb{R}^{B\times N\times 2}\). Node states are initialized by linear projection \(h_i^{(0)}=W_p v_i\).
For each message-passing step \(t\): \(m_i^{(t)}=\phi(W_m h_i^{(t)})\), global context \(g^{(t)}=\frac{1}{N}\sum_i m_i^{(t)}\), update \(h_i^{(t+1)}=\phi\!\left(W_u[h_i^{(t)};g^{(t)}]\right)\), with \(\phi=\mathrm{SiLU}\).
After \(T\) steps (\(T=2\) by default), node states are pooled to \(z_{\text{gnn}}\in\mathbb{R}^{B\times d_{\text{gnn}}}\) (\(d_{\text{gnn}}=64\) default).
Integration: \(z_{\text{gnn}}\) is concatenated with backbone and other branch embeddings.

\textbf{Dynamic head fusion.}
Motivation: enable strict ablations while keeping one unified model definition.
Technical description: the pre-classifier vector is assembled as
\[
z_{\text{all}}=[z_{\text{cnn}};L;z_{\text{att}};z_{\text{gnn}};z_{\text{sobel}}],
\]
where only active components are included, and the FC input dimension is computed programmatically as
\[
d_{\text{fc}}=d_{\text{cnn}}+\mathbb{1}_{L}d_L+\mathbb{1}_{\text{att}}d_{\text{att}}+\mathbb{1}_{\text{gnn}}d_{\text{gnn}}+\mathbb{1}_{\text{sobel}}d_{\text{sobel}}.
\]
For missing optional inputs (e.g., absent landmarks), zero tensors with matching shapes are injected to keep forward passes deterministic.
Integration: branch toggles alter only \(z_{\text{all}}\) composition; training loop, optimizer, and loss interface remain unchanged.

\subsection{B) Separate landmark-detector module}

A dedicated landmark regressor is trained as an independent model.
It uses direct coordinate regression ($2\times$number of keypoints), affine+noise augmentation, and standard regression optimization.
Its purpose is geometric keypoint estimation as a separate task module.

\subsection{C) Separate face-cropper module}

A dedicated face-cropper is trained as a separate multi-task model with two outputs: face confidence and normalized bounding box coordinates.
It supports alternative backbones, boundary-aware bbox loss variants, and geometric/photometric augmentation.
Its purpose is robust face localization as an upstream module.

\subsection{D) Separate K-Means module (creative first idea)}

An original idea was pure non-neural classification via K-Means as a baseline module.
The concept is to use fixed feature spaces (flattened image features, optionally extended with landmarks), cluster into six groups, and map clusters to emotion labels.
